% Figure: a two-transistor current mirror drawn with plain TikZ paths (no
% circuitikz). A diode-connected reference transistor sets a reference current
% through a programming resistor; the output transistor, sharing the same
% base-emitter voltage, copies that current into the load. Schematic level,
% matching the worked example on the constant-current LED driver.
\begin{figure}[htbp]
\centering
\begin{tikzpicture}[
  line width=0.8pt,
  font=\scriptsize,
]
% ===== supply rail =====
\draw (0,4.4) -- (5.2,4.4);
\node[anchor=south] at (2.6,4.45) {$+V_{CC}$};
\filldraw (0,4.4) circle [radius=1.2pt];
\filldraw (5.2,4.4) circle [radius=1.2pt];

% ===== reference branch (left): programming resistor R_ref =====
\draw (0,4.4) -- (0,3.6);
\draw (-0.18,3.6) rectangle (0.18,2.6);
\node[anchor=east] at (-0.28,3.1) {$R_{\text{ref}}$};
\draw (0,2.6) -- (0,2.2);
\filldraw (0,2.2) circle [radius=1.2pt];
\node[engorange, anchor=east] at (-0.28,2.2) {$I_{\text{ref}}$};

% ===== reference transistor Q1 (diode-connected) =====
% collector-base tie
\draw (0,2.2) -- (0,1.9);
% simple bipolar glyph: vertical bar (base) with collector/emitter stubs
\draw[line width=1.4pt] (-0.05,1.9) -- (-0.05,1.1);   % base bar
\draw (0,1.9) -- (-0.55,1.9);                          % collector wire up-left join
\draw (-0.55,1.9) -- (-0.55,1.55) -- (-0.05,1.5);      % collector to base region
\draw (-0.05,1.5) -- (0,1.15);                         % emitter-ish (schematic)
\node[anchor=west] at (0.1,1.5) {$Q_1$};
% base rail shared
\draw (-0.05,1.5) -- (2.6,1.5);
\filldraw (-0.05,1.5) circle [radius=1.2pt];
% emitter of Q1 to ground
\draw (-0.3,1.1) -- (0.2,1.1);
\draw (-0.05,1.1) -- (-0.05,0.4);

% ===== output transistor Q2 =====
\draw[line width=1.4pt] (2.55,1.9) -- (2.55,1.1);      % base bar
\draw (2.6,1.5) -- (2.55,1.5);
\filldraw (2.6,1.5) circle [radius=1.2pt];
\draw (2.55,1.9) -- (2.9,2.2);                          % collector stub
\draw (2.55,1.1) -- (2.9,0.9);                          % emitter stub
\node[anchor=west] at (2.95,1.5) {$Q_2$};
% collector up to load
\draw (2.9,2.2) -- (2.9,2.6);
\filldraw (2.9,2.6) circle [radius=1.2pt];
\node[engorange, anchor=west] at (3.0,2.9) {$I_{\text{out}}\approx I_{\text{ref}}$};

% ===== load (LED symbol as a triangle + bar) at output =====
\draw (5.2,4.4) -- (5.2,3.3);
\draw (5.0,3.3) -- (5.4,3.3);                           % LED bar (cathode)
\draw (5.2,3.3) -- (5.0,3.05) -- (5.4,3.05) -- cycle;   % LED triangle (anode down)
\node[engblue, anchor=west] at (5.5,3.15) {load / LED};
\draw (5.2,3.05) -- (5.2,2.6);
\draw (2.9,2.6) -- (5.2,2.6);

% ===== emitters of Q2 to ground =====
\draw (2.9,0.9) -- (2.9,0.4);

% ===== ground rail =====
\draw (-0.05,0.4) -- (2.9,0.4);
\draw (1.3,0.4) -- (1.3,0.1);
\draw (1.1,0.1) -- (1.5,0.1);
\draw (1.17,-0.02) -- (1.43,-0.02);
\draw (1.25,-0.14) -- (1.35,-0.14);

\end{tikzpicture}
\caption[Current mirror]{A current mirror. The diode-connected reference
transistor $Q_1$ conducts a reference current $I_{\text{ref}}$ set by the
programming resistor, and develops across its base-emitter junction the exact
voltage that a matched transistor needs to carry the same current. The output
transistor $Q_2$ shares that base-emitter voltage, so it copies the reference
current into the load almost independently of the load voltage, the property
that makes the mirror a compact constant-current source for driving a
light-emitting diode or biasing an amplifier stage. The copy is exact only to
the extent the two transistors match and their output resistance is high;
device mismatch and the Early effect set the accuracy.}
\label{fig:vol04ch09:current-mirror}
\end{figure}
